%%%%%%%%%%%%%%%%%%%%%%%%%%%%%%%%%%%%%%%%%%%%%%%%%%%%%%%%%%%%%%%%%%%%%%%%%%%%%%%%%%
\begin{frame}[fragile]\frametitle{}
\begin{center}
{\Large Conclusions}
\end{center}
\end{frame}

%%%%%%%%%%%%%%%%%%%%%%%%%%%%%%%%%%%%%%%%%%%%%%%%%%%%%%%%%%%
\begin{frame}[fragile]\frametitle{Best Practices and Considerations}
      \begin{itemize}
        \item Keep state models simple and clear with necessary information only
        \item Maintain single responsibility in node functions
        % [upgrade-deck: commented out 2026-10-19 | Task 1 | contradicted "return only what changed"]
        % \item Handle exceptions and return new state objects
        \item Handle exceptions inside the node; return only the fields that changed
        \item Use clear conditional logic in edge design
        \item Bound every loop with a step limit; keep cycles simple
        \item Add error handling at critical nodes with fallback mechanisms
        \item Consider performance impacts of checkpoint mechanism
        % [upgrade-deck: commented out 2026-10-19 | Task 1 | vague filler]
        % \item Choose appropriate data processing methods for use cases
      \end{itemize}
\end{frame}

%%%%%%%%%%%%%%%%%%%%%%%%%%%%%%%%%%%%%%%%%%%%%%%%%%%%%%%%%%%
\begin{frame}[fragile]\frametitle{Limitations and What to Watch}
% [upgrade-deck: commented out 2026-10-14 | Task 1 | streaming claims were vague and dated; rewritten for LangGraph 1.0]
% \item Current streaming has long waiting times for LLM nodes
% \item Ideal: Node-level streaming within graph streaming
% \item Market agents provide better step-by-step streaming
% \item LangGraph rapidly evolving with frequent updates
% \item Pre-built components may change in future versions
% \item Monitor documentation for latest features and changes
% \item Core design philosophy remains valuable for learning
% \item Expected improvements in streaming processing capabilities
      \begin{itemize}
        % [upgrade-deck: commented out 2026-10-19 | Task 1 | deprecation claim not confirmed in the official docs]
        % \item The API moves fast: \texttt{create\_react\_agent} became \texttt{create\_agent} in LangChain 1.0.
        \item The API moves fast: LangChain 1.x offers \texttt{create\_agent}, built on LangGraph, for the common agent loop.
        \item Pin versions and re-check the docs before upgrading.
        \item The core ideas (state, nodes, edges, checkpoints) stay stable.
        \item More autonomy means more cost, latency and risk: add step limits and approvals.
        \item An LLM node still waits for the full model call; use token streaming for responsiveness.
      \end{itemize}
\end{frame}

%%%%%%%%%%%%%%%%%%%%%%%%%%%%%%%%%%%%%%%%%%%%%%%%%%%%%%%%%%%
\begin{frame}[fragile]\frametitle{Conclusion}
      \begin{itemize}
        \item State is the key idea: one validated Pydantic context object flows through every node
        \item Nodes are plain functions; edges, including conditional ones, decide the flow
        % [upgrade-deck: commented out 2026-10-19 | Task 1 | LCEL and AgentExecutor are no longer explained in the deck]
        % \item LangGraph addresses LCEL and AgentExecutor limitations effectively
        \item Loops, branches and approvals are drawn in the graph, so the path is visible and testable
        \item Graph-based approach provides intuitive workflow representation
        % [upgrade-deck: commented out 2026-10-19 | Task 1 | vague, merged into the checkpointing bullet below]
        % \item Advanced features support complex AI application development
        % \item State management and persistence enable long-running conversations
        % \item Human-in-the-loop capabilities enhance decision quality
        \item Checkpointing, human approval and streaming make long-running, supervised runs practical
        % [upgrade-deck: commented out 2026-10-19 | Task 2 | the ReAct demo that introduced these helpers is out of the seminar]
        % \item Prebuilt helpers (\texttt{ToolNode}, \texttt{create\_agent}) cover the common ReAct loop
        \item Prebuilt helpers such as \texttt{create\_agent} cover the common agent loop
        % [upgrade-deck: commented out 2026-10-19 | Task 1 | streaming merged into the checkpointing bullet; closing line was filler]
        % \item Streaming responses provide real-time user feedback
        % \item Continuous evolution promises enhanced capabilities for AI development
      \end{itemize}
\end{frame}
